\begin{afterword}
\summaryhead

The post separates truth from confidence. A physical law such as the lightspeed limit may hold always and without exception, and in whatever sense 2 + 2 = 4 is true, it is exactly true. Our confidence in either is a separate matter, as with a coin that shows heads while we have not looked at it. The author is not absolutely confident that 2 + 2 = 4, because the evidence for it could have been hallucinated or misremembered.

The right aim, the post says, is calibration: statements held at 99\% should come true 99 times in 100. People who say they are 99\% sure are less accurate than that; the author points to a book chapter for the experiments. To claim 99.9999\% is to claim that you could make a million statements, about a year of talking, with about one error. Higher levels need a trillion statements or more. The post closes with Rafal Smigrodski's remark that a probability of 1 could never be revised, whatever one learned.

\responsehead

The post's two main points are sound. That a statement can be exactly true while our confidence in it falls short is the standard distinction between truth and belief, and the post states it clearly with the unseen coin. The closing argument is also correct: under Bayes's rule a probability of 1 cannot be lowered by any evidence. The one empirical claim, that people who say they are 99\% sure are right less often than that, is supported by the chapter the post cites, which reports that answers given odds of 100 to 1 were right 73\% of the time.

Much of the post restates ``Absolute Authority,'' published the day before. The million statements in a year of talking appear there almost word for word, and so does the point that a belief one could revise cannot have probability 1. The new illustration contains a slip. At the pace the post gives, a trillion statements take about 950,000 years, not ``a hundred human lifetimes.'' The error makes the post's own case look weaker than it is. The claim that a confidence of $1 - 1/(3 \uparrow\uparrow\uparrow 3)$ is not much closer to 1 than 90\% also depends on a scale the post does not give; the log-odds scale arrives the next day.

The post leaves out one technical point that matters for what follows. Standard probability theory gives probability 1 to every logical truth, and that SS0 + SS0 = SSSS0 follows from the Peano axioms is such a truth. Holding it at less than 1, as the post does, is a known departure from the standard theory, argued for by Hacking and Garber among others. The post takes that side without saying so. The next day's post proposes to change probability theory itself; a reader of this one has not been told that the theory, as it stands, requires certainty in logical truths.

In short: a sound distinction between truth and confidence and a correct point about updating, largely repeated from the previous day's post, with no mention that standard probability theory gives logical truths probability 1.
\end{afterword}
